\chapter*{GLOSARIO}
\addcontentsline{toc}{chapter}{Glosario}

\begin{description}
  \item[Arquitectura de plugins:] organización de un sistema en un núcleo estable y módulos intercambiables que implementan un contrato común y se pueden agregar o reemplazar sin modificar el núcleo.
  \item[Capa de abstracción de hardware (HAL):] componente que oculta los detalles de cada dispositivo detrás de una interfaz uniforme, de modo que un dispositivo real puede reemplazarse por uno simulado.
  \item[Clic por permanencia (\textit{dwell}):] activación de un control cuando la mirada permanece dentro de un radio durante un tiempo mínimo.
  \item[Consistencia bidireccional:] en EDDIE, mantener equivalentes el documento impreso y su versión PDF, de modo que una marca hecha en uno se refleje en el otro.
  \item[Contrato de interfaz:] conjunto de operaciones que un módulo se compromete a implementar; en este trabajo, una clase abstracta de Python.
  \item[Deuda técnica:] costo futuro de retrabajo derivado de elegir soluciones rápidas en lugar de un diseño mantenible.
  \item[Eye tracker:] dispositivo que estima, a una frecuencia fija, el punto al que mira una persona.
  \item[Fijación:] período en que la mirada permanece estable sobre una región.
  \item[Homografía:] transformación proyectiva entre dos planos, usada para convertir coordenadas de la cámara en coordenadas de la proyección.
  \item[Interacción humano-información (HII):] campo que estudia cómo las personas buscan, evalúan y usan información en distintos formatos.
  \item[Latencia:] tiempo entre una acción y su resultado; en este trabajo, el tiempo que el orquestador tarda en ejecutar un plugin o un ciclo completo.
  \item[MECIVA PS:] metodología para la evaluación y validación de productos de software desarrollados para EDDIE \citep{ibaceta2023}.
  \item[Modelo de lenguaje (LLM):] modelo de aprendizaje automático que genera texto a partir de una instrucción; en EDDIE responde los servicios de búsqueda.
  \item[OCR:] reconocimiento óptico de caracteres; conversión de la imagen de un texto en caracteres.
  \item[Orquestador:] componente central que controla el ciclo de vida de los plugins y coordina el flujo de datos entre ellos.
  \item[Percentil 95 (p95):] valor bajo el cual queda el 95\% de las mediciones.
  \item[Prototipo desechable:] pieza de software acotada que se construye para evaluar la factibilidad de una función y luego se descarta.
  \item[QuadPoints:] coordenadas de los cuatro vértices de cada fragmento de texto en un PDF, usadas para resaltarlo.
  \item[Realidad aumentada espacial:] superposición de información digital directamente sobre los objetos físicos mediante un proyector.
  \item[\textit{Science gateway}:] plataforma que integra instrumentos, datos y flujos de trabajo científicos.
  \item[Software científico:] software desarrollado \textit{ad hoc} para soportar una tarea de investigación.
  \item[SonarQube:] plataforma de análisis estático de la calidad, la seguridad y la mantenibilidad del código.
  \item[Toque de Midas:] activación involuntaria de controles por el solo hecho de mirarlos.
\end{description}
